\fancyhead[L]{\small ANZEIGEN.}
\SourceLeaf{346}
\pdfbookmark[0]{Anzeigen. Disquisitiones generales circa superficies curvas.}{anzeigen-curvas}
\begin{center}
\Large ANZEIGEN.\par\medskip
\normalsize\rule{.85\textwidth}{.3pt}\par
Göttingische gelehrte Anzeigen.\quad 5. November 1827.\par
\rule{.85\textwidth}{.3pt}
\end{center}
\hspace*{\parindent}Am 8. Oktober überreichte Herr Hofrat \textsc{Gauss} der Königlichen Sozietät einen Vortrag.
\begin{center}\textit{Disquisitiones generales circa superficies curvas.}\end{center}
\hspace*{\parindent}Die Geometer haben sich eingehend mit allgemeinen Untersuchungen über gekrümmte Flächen beschäftigt, und ihre Ergebnisse machen einen bedeutenden Teil der höheren Geometrie aus. Doch ist der Gegenstand noch so weit davon entfernt, erschöpft zu sein, dass man vielmehr sagen kann, bisher sei erst ein kleiner Teil eines höchst fruchtbaren Feldes bebaut worden. Schon vor einigen Jahren hat der Verfasser versucht, diesem Gebiet eine neue Seite abzugewinnen, indem er die Aufgabe löste, alle Darstellungen einer gegebenen Fläche auf einer anderen zu finden, bei denen die kleinsten Teile ähnlich bleiben. Die vorliegende Abhandlung soll weitere neue Gesichtspunkte eröffnen und einen Teil der dadurch zugänglichen neuen Wahrheiten entwickeln. Wir werden hier darüber berichten, soweit sich dies ohne zu große Weitschweifigkeit verständlich machen lässt. Vorab müssen wir jedoch darauf hinweisen, dass sowohl die neuen Begriffsbildungen als auch die Lehrsätze teilweise noch Einschränkungen oder genauere Bestimmungen erfordern, wenn die größtmögliche Allgemeinheit erfasst werden soll. Diese müssen hier übergangen werden.
